\section{Introduction}
\label{sec:intro}

A tool call is an executive act. When an LLM agent reads a file, writes to a database, sends an email, executes a shell command, calls a third-party API, or pushes a commit, it takes an action on behalf of a user, against an external system, with consequences in the world. The agent receives delegated authority through a system prompt or a developer-issued credential; it dispatches that authority against a resource it does not own; it produces an effect that, if wrong, must be undone.

The agentic-AI safety literature has converged on a particular framing of this problem. Safety is treated as a property of the model: how to make the model want the right thing, refuse the wrong thing, and behave consistently between evaluation and deployment. Constitutional AI~\cite{baiConstitutionalAI2022} shapes preferences at training time. Reinforcement learning from human feedback~\cite{ouyangRLHF2022} encodes operator judgement into reward. Scalable oversight~\cite{leikeScalableOversight2018,bowmanScalableOversight2022} and debate~\cite{christianoDebate2018} aim to extend supervision past the operator's direct capability. Alignment-faking~\cite{hubingerAlignmentFaking2024} and situational-awareness~\cite{berglundSituationalAwareness2023} research sharpen the threat model: a model may behave well under evaluation and differently in deployment, and the cause may be the model's own situational reasoning rather than an outer optimisation failure. Each thread treats the model as the object of safety analysis.

We argue that safety is also, separately, a property of the substrate that admits the tool call. The substrate may decline to execute a tool call that does not carry a typed proof of reversibility, regardless of what the model that issued the call wants, regardless of whether the model successfully evaded its training signal, and regardless of whether the operator was manipulated into approving it. The grammar of admission is independent of the grammar of generation.

The structural argument is short. A tool call has the shape of a delegated executive act: a typed envelope, an authorisation chain, a target, an effect, and (in the destructive case) a need for reversal. A substrate that requires every destructive tool call to be constructed only with a positive time-to-live (TTL) witness and a typed rollback receipt slot refuses tool calls that the model cannot, by construction, prove are reversible. The model may be misaligned; it cannot forge a rollback witness for an action whose effects it cannot undo. The substrate may be type-checked external to the model.

This paper develops the argument in five steps. We name the existing alignment-research stack and what it currently does not cover (Section~\ref{sec:background}). We argue that a tool call is structurally an executive act, with the same shape as an emergency response action in an endpoint-security system or a temporally-bounded administrative order in a legal system (Section~\ref{sec:executive-acts}). We sketch the formal grammar of admission and cite the structural theorems already discharged in companion work on the underlying substrate (Section~\ref{sec:grammar}). We describe an implementation instance (Section~\ref{sec:impl}), state the threat model that a substrate-layer guarantee survives (Section~\ref{sec:threat}), and discuss how training-layer and admission-layer safety compose (Section~\ref{sec:discussion}). We close with limitations (Section~\ref{sec:limitations}) and a conclusion (Section~\ref{sec:conclusion}).

The contribution is structural rather than empirical. The paper does not introduce a new training-time technique, an interpretability primitive, or an evaluation benchmark. It proposes that two safety layers exist, that they compose, and that the runtime-substrate layer holds under adversarial assumptions about the model that training-time interventions cannot match. Neither layer alone is sufficient; the combination is the construction the literature should be exploring.
